\noindent\hspace*{\parindent}Every summer, cities face the challenge of keeping people cool in open-air and semi-open spaces, ranging from outdoor bus shelters and rail platforms to campus courtyards and covered airport curbsides. Unlike enclosed buildings, these spaces constantly exchange air with their surroundings, allowing some of the cooling to escape before it reaches the people it is meant to help. Finding an effective cooling strategy therefore requires balancing the trade-offs between energy consumption, human comfort, and financial and resource constraints. To address this challenge, we developed a mathematical framework to compare cooling systems based on their efficiency, effectiveness, and cost.

We began by examining how cooling systems interact with their surroundings. Using psychrometric relationships and steady-state heat and moisture balances, we estimated how much cooling actually reaches occupants after mixing with outdoor air. For rail platforms, we used the same balances with higher outdoor air exchange and baseline airspeed to represent their more open design. However, a reduction in air temperature does not always translate to an equivalent improvement in human comfort. To account for this, we applied the Gagge two-node model of human thermoregulation, which describes heat exchange between the body’s core, skin, and surrounding environment. We then calculated Standard Effective Temperature (SET) to measure the comfort experienced by occupants, and defined cooling efficiency as the SET reduction per kWh of electricity.

After measuring the comfort each cooling system provides, we needed to determine which option would be most effective under different conditions. Our decision model can screen options against budget, electricity, and water limits, and applies minimax regret to select the system with the smallest worst-case extra discomfort relative to the best feasible option across different weather scenarios. Finally, we conducted deterministic sensitivity analysis to examine how changes in our assumptions affected the results and whether they led to different cooling recommendations.

We applied our framework to 12 venues across Dallas, Phoenix, and Boston under peak-heat, hot-humid, and hot-dry conditions. With the assumed passenger airspeeds, the Fan delivered the most cooling per kWh in all 36 venue–weather cases. With no budget or resource limits applied, the larger evaporative cooler had the lowest hot-discomfort loss and was the minimax choice at all 12 venues. At the Boston bus shelter, however, the fan provided most of the larger cooler’s peak-hour relief at a tenth of the electricity. We also found that higher outdoor air exchange reduced the benefit of evaporative cooling, and that small changes in the hot-humid outdoor temperature or the upper comfort threshold switched Boston’s recommendation between the fan and the larger cooler.

Our findings show that the cooling system providing the most comfort is not always the most efficient choice. Operators should consider local weather conditions, airflow, and available resources when deciding which system to use. Fans may be sufficient when only moderate cooling is needed, while evaporative coolers can provide greater relief during peak heat. Even the larger cooler left every venue above the comfort target during peak heat, so additional measures such as shade and access to air-conditioned shelter may be necessary to keep public spaces comfortable and safe.
